Embodied work includes any fulfillment step that depends on physical presence, place-specific context, local access, real-world timing, or human witnessing. Common examples include onsite visits, field inspection, pickup or dropoff, event attendance, device or environment interaction, in-person verification, and handoff that requires a signature or observed confirmation. These steps differ from pure digital tasks because they cannot be reduced to generation quality alone. Their correctness depends on where, when, and by whom the work occurred.

For that reason, an orchestration model should represent embodied work explicitly rather than treating it as a weak annotation on a generic task. At minimum, the system should be able to distinguish among execution modalities such as remote digital work, remote synchronous interaction, onsite visit, field inspection, pickup or dropoff, human review, and witnessed handoff. It should also be able to represent planning constraints such as geography, time window, access dependency, safety requirement, or required proof type.

These requirements do not necessarily imply a new root object. They can be represented through canonical request fields, structured matching criteria, fulfillment-step metadata, or proof requirements attached to the execution lane. The crucial point is that the information must be durable enough to affect routing and validation. If the need for onsite presence exists only in one freeform paragraph, it is easy for a planner to ignore it and still appear productive.

Explicit embodied representation also changes how supply is modeled. A supply should not be described only by abstract capability labels. It should be possible to express whether a supply can provide local presence, what geography it can serve, whether it can perform inspections or witness tasks, whether it can handle access-constrained sites, and what proof surfaces it can produce. This lets the system distinguish a strong digital analyst from a field-capable operator even if both could generate polished reports.

Finally, explicit representation changes completion rules. A summary that says ``inspection completed'' is not the same as an inspection artifact. A generated checklist is not the same as a geostamped site photo, a witness signature, or a time-bound access log. The model should therefore require proof appropriate to the execution mode. Otherwise, the system cannot distinguish real completion from narrative closure.
